\documentclass[11pt]{article}

\usepackage{amsmath,amssymb,amsthm}

\newtheorem{proposition}{Proposition}
\newtheorem{lemma}[proposition]{Lemma}
\newtheorem{corollary}[proposition]{Corollary}
\newcommand{\E}{\mathbb{E}}
\newcommand{\calI}{\mathcal{I}}
\newcommand{\ind}{\mathbf{1}}
\newcommand{\norm}[1]{\lVert #1\rVert}
\newcommand{\R}{\mathbb{R}}
\DeclareMathOperator{\tr}{tr}
\DeclareMathOperator{\Var}{Var}

\begin{document}

\section{The SYK model and its partition function}
\label{sec:syk-model}

Fix integers $N,q\in\mathbb{N}$ such that $0<q\leq N$ and $q$ is even. Write
\begin{equation}
  \mathcal{I}_{N,q}
  =
  \{I\subseteq\{1,\ldots,N\}: |I|=q\}.
\end{equation}
Thus, an interaction label is a $q$-element subset of the Majorana labels.

Let $\mathcal{A}$ be a complex star algebra with generators
$\psi_1,\ldots,\psi_N\in\mathcal{A}$ satisfying the Majorana relations
\begin{equation}
  \psi_i\psi_j+\psi_j\psi_i=2\delta_{ij},
  \qquad
  \psi_i^*=\psi_i.
\end{equation}
For $I=\{i_1<\cdots<i_q\}\in\mathcal{I}_{N,q}$, define the ordered
Majorana monomial
\begin{equation}
  \psi_I=\psi_{i_1}\cdots\psi_{i_q}.
\end{equation}
Reversing the order of the product gives
\begin{equation}
  \psi_I^*=(-1)^{q/2}\psi_I,
\end{equation}
where the sign is understood as a scalar in $\mathbb{C}$. Consequently, the
phase
\begin{equation}
  \phi_q=i^{q/2}
\end{equation}
produces a self-adjoint Hamiltonian from arbitrary real couplings
$J=(J_I)_{I\in\mathcal{I}_{N,q}}$:
\begin{equation}
  H(J)
  =
  \sum_{I\in\mathcal{I}_{N,q}} J_I\phi_q\psi_I,
  \qquad
  H(J)^*=H(J).
  \label{eq:syk-hamiltonian}
\end{equation}

For the thermodynamic quantities, take a finite-dimensional matrix
representation $H\in\operatorname{Mat}_d(\mathbb{C})$ of
\eqref{eq:syk-hamiltonian}. For inverse temperature $\beta\geq0$, set
\begin{align}
  \Gamma_\beta(H) &= \exp(-\beta H), \\
  Z(\beta;H) &= \operatorname{Tr}\bigl(\Gamma_\beta(H)\bigr), \\
  \langle O\rangle_{\beta,H}
    &= \frac{\operatorname{Tr}\bigl(O\Gamma_\beta(H)\bigr)}
            {Z(\beta;H)}.
\end{align}
The trace is unnormalized, so $\Gamma_0(H)=I_d$ and $Z(0;H)=d$.

\section{Gaussian disorder and two replicas}
\label{sec:two-replicas}

To obtain the random SYK model, let $(g_I)_{I\in\mathcal{I}_{N,q}}$ be
independent standard real Gaussian variables and choose a deterministic
scale $\sigma_N\geq0$. We use the couplings
\begin{equation}
  J_I=\sigma_N g_I,
  \qquad
  Z_N(\beta;g)=Z\bigl(\beta;H(J)\bigr),
  \qquad
  F_N(\beta)=\mathbb{E}\bigl[Z_N(\beta;g)\bigr].
\end{equation}

Let $\widetilde g$ be an independent copy of $g$. For
$\rho\in[-1,1]$, define the correlated disorder
\begin{equation}
  g_I^{(\rho)}
  =
  \rho g_I+\sqrt{1-\rho^2}\,\widetilde g_I,
  \qquad I\in\mathcal{I}_{N,q},
\end{equation}
and the two-replica partition function
\begin{equation}
  \mathcal{G}_N(\rho,\beta)
  =
  \mathbb{E}\left[
    Z_N(\beta;g)Z_N\bigl(\beta;g^{(\rho)}\bigr)
  \right].
\end{equation}

\begin{proposition}[Mehler decomposition]
  Let $P_r$ denote orthogonal projection onto the $r$th Gaussian Wiener
  chaos. Then
  \begin{equation}
    \mathcal{G}_N(\rho,\beta)
    =
    F_N(\beta)^2
    +
    \sum_{r\geq1}\rho^r
      \left\lVert P_r Z_N(\beta;g)\right\rVert_{L^2}^2.
  \end{equation}
  In particular,
  \begin{equation}
    \mathbb{E}\left[Z_N(\beta;g)^2\right]
    =
    F_N(\beta)^2
    +
    \sum_{r\geq1}
      \left\lVert P_r Z_N(\beta;g)\right\rVert_{L^2}^2.
  \end{equation}
\end{proposition}

% BEGIN INSERTED SECOND-MOMENT SECTION
\section{A correlated two-replica calculation for the partition function}
\label{sec:partition-function-second-moment}

Throughout this section, let
\begin{equation}\label{eq:annealed-partition-function}
  F_N(\beta)=\E Z_N(\beta).
\end{equation}
Let $\widetilde g=(\widetilde g_I)_{I\in\calI_N}$ be an independent copy of
$g=(g_I)_{I\in\calI_N}$.  For $\rho\in[-1,1]$, define
\begin{equation}\label{eq:rho-correlated-disorder}
  g_I^{(\rho)}
  =
  \rho g_I+\sqrt{1-\rho^2}\,\widetilde g_I,
  \qquad I\in\calI_N,
\end{equation}
and set
\begin{equation}\label{eq:rho-correlated-second-moment}
  \mathcal G_N(\rho,\beta)
  =
  \E\left[
    Z_N(\beta;g)Z_N(\beta;g^{(\rho)})
  \right].
\end{equation}

\begin{proposition}[Mehler decomposition]\label{prop:Zn-Mehler}
  For every $\rho\in[-1,1]$ and $\beta\in\R$,
  \begin{equation}\label{eq:Zn-Mehler}
    \mathcal G_N(\rho,\beta)
    =
    F_N(\beta)^2
    +
    \sum_{r\geq1}\rho^r\norm{P_rZ_N(\beta)}_2^2.
  \end{equation}
  In particular,
  \begin{equation}\label{eq:Zn-second-moment-chaos}
    \E[Z_N(\beta)^2]
    =
    F_N(\beta)^2
    +
    \sum_{r\geq1}\norm{P_rZ_N(\beta)}_2^2.
  \end{equation}
\end{proposition}

\begin{proof}
  Write the Hermite expansion as
  \begin{equation}\label{eq:Zn-Hermite-expansion}
    Z_N(\beta;g)
    =
    \sum_{\alpha}z_{N,\alpha}(\beta)H_\alpha(g).
  \end{equation}
  The correlated Gaussian identity gives
  \begin{equation}\label{eq:correlated-Hermite-identity}
    \E\left[
      H_\alpha(g)H_\gamma(g^{(\rho)})
    \right]
    =
    \ind_{\{\alpha=\gamma\}}\rho^{|\alpha|}\alpha!.
  \end{equation}
  Substituting \eqref{eq:Zn-Hermite-expansion} into
  \eqref{eq:rho-correlated-second-moment} and grouping terms according to
  $|\alpha|$ proves \eqref{eq:Zn-Mehler}.  Taking $\rho=1$ proves
  \eqref{eq:Zn-second-moment-chaos}.
\end{proof}

\begin{lemma}[Exact first and second chaoses of $Z_N$]
  \label{lem:Zn-first-second-chaos}
  The first chaos vanishes:
  \begin{equation}\label{eq:Zn-first-chaos-zero}
    P_1Z_N(\beta)=0.
  \end{equation}
  Moreover, there exists a scalar $b_N(\beta)$ such that
  \begin{equation}\label{eq:Zn-second-chaos}
    P_2Z_N(\beta)
    =
    b_N(\beta)
    \sum_{I\in\calI_N}(g_I^2-1),
  \end{equation}
  and
  \begin{equation}\label{eq:bn-exact}
    2M_Nb_N(\beta)
    =
    \beta F_N'(\beta).
  \end{equation}
  Consequently,
  \begin{equation}\label{eq:Zn-second-chaos-norm}
    \norm{P_2Z_N(\beta)}_2^2
    =
    \frac{\beta^2}{2M_N}F_N'(\beta)^2.
  \end{equation}
\end{lemma}

\begin{proof}
  The same Majorana sign-gauge transformation used for $L_N$ gives
  \begin{equation}\label{eq:Zn-gauge-invariance}
    Z_N(\beta,g)=Z_N(\beta,\chi_sg),
    \qquad s\in\{\pm1\}^N.
  \end{equation}
  The corresponding Hermite parity rule implies
  \eqref{eq:Zn-first-chaos-zero} and eliminates all off-diagonal terms
  $g_Ig_J$ with $I\neq J$ from the second chaos.  Permutation symmetry then
  gives \eqref{eq:Zn-second-chaos}.

  By the Hermite coefficient formula,
  \begin{equation}\label{eq:bn-second-derivative}
    b_N(\beta)
    =
    \frac12\E[\partial_{g_I}^2Z_N(\beta)].
  \end{equation}
  Summing over $I$ and applying Gaussian integration by parts yields
  \begin{equation}\label{eq:bn-gibp}
    \begin{aligned}
      2M_Nb_N(\beta)
      &=
      \sum_{I\in\calI_N}\E[\partial_{g_I}^2Z_N(\beta)]
      \\
      &=
      \sum_{I\in\calI_N}\E[g_I\partial_{g_I}Z_N(\beta)].
    \end{aligned}
  \end{equation}
  Since $H_N(g)$ is linear in $g$,
  \begin{equation}\label{eq:Zn-radial-derivative}
    \sum_{I\in\calI_N}g_I\partial_{g_I}Z_N(\beta)
    =
    \beta\partial_\beta Z_N(\beta).
  \end{equation}
  Taking expectation proves \eqref{eq:bn-exact}.  Finally,
  $\Var(g_I^2-1)=2$, so \eqref{eq:Zn-second-chaos-norm} follows from
  \eqref{eq:Zn-second-chaos} and \eqref{eq:bn-exact}.
\end{proof}

\begin{corollary}[Exact second-moment reduction]
  \label{cor:Zn-exact-second-moment-reduction}
  For every $\beta\in\R$,
  \begin{equation}\label{eq:Zn-exact-second-moment-reduction}
    \frac{\E[Z_N(\beta)^2]}{F_N(\beta)^2}
    =
    1
    +
    \frac{\beta^2}{2M_N}
    \left(
      \partial_\beta\log F_N(\beta)
    \right)^2
    +
    \mathcal R_{N,\geq3}^{Z}(\beta),
  \end{equation}
  where
  \begin{equation}\label{eq:Zn-higher-chaos-remainder}
    \mathcal R_{N,\geq3}^{Z}(\beta)
    =
    \frac1{F_N(\beta)^2}
    \sum_{r\geq3}\norm{P_rZ_N(\beta)}_2^2
    \geq0.
  \end{equation}
  More generally,
  \begin{equation}\label{eq:Zn-rho-exact-reduction}
    \frac{\mathcal G_N(\rho,\beta)}{F_N(\beta)^2}
    =
    1
    +
    \rho^2\frac{\beta^2}{2M_N}
    \left(
      \partial_\beta\log F_N(\beta)
    \right)^2
    +
    \sum_{r\geq3}\rho^r
    \frac{\norm{P_rZ_N(\beta)}_2^2}{F_N(\beta)^2}.
  \end{equation}
\end{corollary}

\begin{proof}
  Combine Proposition~\ref{prop:Zn-Mehler} with
  Lemma~\ref{lem:Zn-first-second-chaos}.
\end{proof}

We next compute the first two non-zero Taylor coefficients of
\eqref{eq:Zn-exact-second-moment-reduction}.  Define
\begin{equation}\label{eq:Sn-sign-sum}
  S_N
  =
  \sum_{I,J\in\calI_N}(-1)^{|I\cap J|}
  =
  M_N
  \sum_{j=0}^q
    (-1)^j
    \binom qj
    \binom{N-q}{q-j},
\end{equation}
and
\begin{equation}\label{eq:Un-triples}
  U_N
  =
  \begin{cases}
    \displaystyle
    \frac{N!}{
      6\left(N-\frac{3q}{2}\right)!
      \left(\left(\frac q2\right)!\right)^3
    },
    & N\geq\frac{3q}{2},
    \\
    0,
    & N<\frac{3q}{2}.
  \end{cases}
\end{equation}

\begin{proposition}[Taylor expansion of the normalized second moment]
  \label{prop:Zn2-Taylor-expansion}
  For fixed $N$ and $\beta\to0$,
  \begin{equation}\label{eq:Fn-Taylor-expansion}
    F_N(\beta)
    =
    1
    +
    \frac{\beta^2}{2}\sigma_N^2M_N
    +
    \frac{\beta^4}{24}\sigma_N^4
    \left(
      2M_N^2+S_N
    \right)
    +
    O_{N,q}(\beta^6).
  \end{equation}
  Moreover,
  \begin{equation}\label{eq:Zn2-ratio-Taylor-expansion}
    \begin{aligned}
      \frac{\E[Z_N(\beta)^2]}{F_N(\beta)^2}
      ={}&
      1
      +
      \frac{\beta^4}{2}\sigma_N^4M_N
      \\
      &+
      \beta^6\sigma_N^6
      \left\{
        \frac{S_N-M_N^2}{6}
        +
        \ind_{\{4\mid q\}}U_N
      \right\}
      +
      O_{N,q}(\beta^8).
    \end{aligned}
  \end{equation}
\end{proposition}

\begin{proof}
  The Clifford relations imply
  \begin{equation}\label{eq:Clifford-two-four-traces}
    \tr(\psi_I\psi_J)
    =
    \ind_{\{I=J\}},
    \qquad
    \tr(\psi_I\psi_J\psi_I\psi_J)
    =
    (-1)^{|I\cap J|}.
  \end{equation}
  Hence Wick's formula gives
  \begin{equation}\label{eq:H2-H4-annealed}
    \E\tr(H_N^2)
    =
    \sigma_N^2M_N,
    \qquad
    \E\tr(H_N^4)
    =
    \sigma_N^4
    \left(
      2M_N^2+S_N
    \right).
  \end{equation}
  Expanding the exponential proves \eqref{eq:Fn-Taylor-expansion}.
  It follows that
  \begin{equation}\label{eq:log-Fn-derivative-Taylor}
    \partial_\beta\log F_N(\beta)
    =
    \sigma_N^2M_N\beta
    +
    \frac{\sigma_N^4}{6}
    \left(
      S_N-M_N^2
    \right)\beta^3
    +
    O_{N,q}(\beta^5).
  \end{equation}
  Therefore, the second-chaos term in
  \eqref{eq:Zn-exact-second-moment-reduction} equals
  \begin{equation}\label{eq:Zn-second-chaos-Taylor}
    \begin{aligned}
      \frac{\beta^2}{2M_N}
      \left(
        \partial_\beta\log F_N(\beta)
      \right)^2
      ={}&
      \frac{\beta^4}{2}\sigma_N^4M_N
      \\
      &+
      \frac{\beta^6}{6}\sigma_N^6
      \left(
        S_N-M_N^2
      \right)
      +
      O_{N,q}(\beta^8).
    \end{aligned}
  \end{equation}

  It remains to identify the order-$\beta^6$ contribution of the third chaos.
  Every third-chaos multi-index with a repeated hyperedge label violates the
  Majorana gauge-parity rule.  Thus only products $g_Ig_Jg_K$ with distinct
  $I,J,K\in\calI_N$ can contribute.  The coefficient of such a product in
  the order-$\beta^3$ term of $Z_N(\beta)$ can be non-zero only if
  \begin{equation}\label{eq:triple-symmetric-difference}
    I\mathbin{\triangle}J\mathbin{\triangle}K=\varnothing.
  \end{equation}
  In that case, every vertex belongs to exactly two of $I,J,K$, and hence
  \begin{equation}\label{eq:triple-intersections}
    |I\cap J|
    =
    |I\cap K|
    =
    |J\cap K|
    =
    \frac q2.
  \end{equation}
  The number of unordered triples satisfying
  \eqref{eq:triple-symmetric-difference} is $U_N$.  The six orderings of such
  a triple have the same trace when $4\mid q$, whereas they cancel in pairs
  when $q\equiv2\pmod4$.  Since
  \begin{equation}\label{eq:triple-trace-magnitude}
    \left|
      \tr(\psi_I\psi_J\psi_K)
    \right|
    =
    1,
  \end{equation}
  we obtain
  \begin{equation}\label{eq:P3Zn-leading-norm}
    \norm{P_3Z_N(\beta)}_2^2
    =
    \ind_{\{4\mid q\}}
    \beta^6\sigma_N^6U_N
    +
    O_{N,q}(\beta^8).
  \end{equation}
  The total contribution of chaoses of order at least four is
  $O_{N,q}(\beta^8)$.  Substituting
  \eqref{eq:Zn-second-chaos-Taylor} and
  \eqref{eq:P3Zn-leading-norm} into
  \eqref{eq:Zn-exact-second-moment-reduction} proves
  \eqref{eq:Zn2-ratio-Taylor-expansion}.
\end{proof}

\begin{corollary}[Large-$N$ asymptotics of the first Taylor coefficients]
  \label{cor:Zn2-large-n-Taylor-coefficients}
  As $N\to\infty$,
  \begin{equation}\label{eq:Sn-asymptotic}
    \frac{S_N}{M_N^2}
    =
    1-\frac{2q^2}{N}+O_q(N^{-2}),
  \end{equation}
  and
  \begin{equation}\label{eq:Zn2-beta4-coefficient}
    \frac12\sigma_N^4M_N
    =
    \frac{(q-1)!}{2q}
    N^{2-q}
    \left(
      1+O_q(N^{-1})
    \right).
  \end{equation}
  Furthermore,
  \begin{equation}\label{eq:Zn2-beta6-second-chaos-coefficient}
    \frac16\sigma_N^6
    \left(
      S_N-M_N^2
    \right)
    =
    -\frac{(q-1)!}{3}
    N^{2-q}
    \left(
      1+O_q(N^{-1})
    \right),
  \end{equation}
  whereas
  \begin{equation}\label{eq:Zn2-beta6-third-chaos-coefficient}
    \sigma_N^6U_N
    =
    O_q\left(N^{3-3q/2}\right)
    =
    o\left(N^{2-q}\right).
  \end{equation}
  Consequently,
  \begin{equation}\label{eq:Zn2-fourth-derivative-limit}
    \lim_{N\to\infty}
    N^{q-2}\frac1{4!}
    \left.
      \partial_\beta^4
      \frac{\E[Z_N(\beta)^2]}{F_N(\beta)^2}
    \right|_{\beta=0}
    =
    \frac{(q-1)!}{2q},
  \end{equation}
  and
  \begin{equation}\label{eq:Zn2-sixth-derivative-limit}
    \lim_{N\to\infty}
    N^{q-2}\frac1{6!}
    \left.
      \partial_\beta^6
      \frac{\E[Z_N(\beta)^2]}{F_N(\beta)^2}
    \right|_{\beta=0}
    =
    -\frac{(q-1)!}{3}.
  \end{equation}
\end{corollary}

\begin{proof}
  For fixed $I\in\calI_N$ and uniformly sampled $J\in\calI_N$,
  $|I\cap J|$ has the hypergeometric distribution.  Therefore,
  \begin{equation}\label{eq:sign-hypergeometric}
    \frac{S_N}{M_N^2}
    =
    \E\left[(-1)^{|I\cap J|}\right]
    =
    1-2\mathbb P(|I\cap J|\text{ is odd}).
  \end{equation}
  Since
  \begin{equation}\label{eq:intersection-one-asymptotic}
    \mathbb P(|I\cap J|=1)
    =
    \frac{q\binom{N-q}{q-1}}{\binom Nq}
    =
    \frac{q^2}{N}+O_q(N^{-2}),
  \end{equation}
  and $\mathbb P(|I\cap J|\geq2)=O_q(N^{-2})$,
  \eqref{eq:Sn-asymptotic} follows.  The remaining statements follow from
  $M_N\sim N^q/q!$, $\sigma_N^2=(q-1)!N^{-(q-1)}$, and
  $U_N=O_q(N^{3q/2})$.
\end{proof}

\paragraph{Remaining high-temperature estimate.}
Corollary~\ref{cor:Zn-exact-second-moment-reduction} shows that a CLT for
$Z_N(\beta)$ follows once one proves, for fixed sufficiently small
$\beta\neq0$,
\begin{equation}\label{eq:Zn2-target-second-moment}
  \frac{\E[Z_N(\beta)^2]}{F_N(\beta)^2}
  =
  1
  +
  \frac{\beta^2}{2M_N}
  \left(
    \partial_\beta\log F_N(\beta)
  \right)^2
  +
  o\left(N^{2-q}\right).
\end{equation}
Indeed, all terms in \eqref{eq:Zn-higher-chaos-remainder} are non-negative,
so \eqref{eq:Zn2-target-second-moment} is equivalent to
\begin{equation}\label{eq:Zn-higher-chaos-target}
  \sum_{r\geq3}\norm{P_rZ_N(\beta)}_2^2
  =
  o\left(F_N(\beta)^2N^{2-q}\right).
\end{equation}
The Taylor calculation above verifies the expected second-chaos dominance
through order $\beta^6$: the third-chaos contribution at that order is smaller
by a factor of order $N^{1-q/2}$.
% END INSERTED SECOND-MOMENT SECTION

\begin{thebibliography}{GPSZ26}

\bibitem[GPSZ26]{GPSZ2026}
David Gamarnik, Francisco Pernice, Alexander Schmidhuber, and Alexander
Zlokapa.
\newblock The free energy limit of the {SYK} model at high temperature.
\newblock arXiv:2605.02768, 2026.

\bibitem[ACKK25]{ACKK2025}
Eric R. Anschuetz, Chi-Fang Chen, Bobak T. Kiani, and Robbie King.
\newblock Strongly interacting fermions are non-trivial yet non-glassy.
\newblock Physical Review Letters, 135(3):030602, 2025.

\bibitem[Wil67]{Wilcox1967}
R. M. Wilcox.
\newblock Exponential operators and parameter differentiation in quantum physics.
\newblock Journal of Mathematical Physics, 8(4):962--982, 1967.

\end{thebibliography}
\end{document}
